% === Introduction ===
\section{Introduction}
\label{sec:introduction}

Your introduction goes here. Describe the problem, motivation, and contributions of your work.

% --- Example: how to cite references ---
For example, the Transformer architecture~\cite{vaswani2017attention} has been widely adopted in natural language processing, computer vision, and multimodal learning. Its self-attention mechanism enables modeling long-range dependencies, making it the backbone of many state-of-the-art systems. Deep residual networks~\cite{he2016deep} enabled training of very deep models by introducing skip connections, which alleviate the vanishing gradient problem and have become a standard component in modern architectures. You can also cite multiple works at once~\cite{devlin2019bert,goodfellow2014generative,kingma2014adam}.

Despite these advances, several challenges remain. [Describe the key challenges or gaps in the current literature that motivate your work. Explain why existing methods are insufficient and what specific limitations you aim to address.]

To address these challenges, we propose \method, a [brief one-sentence description of your approach]. The key insight behind our method is [describe your core idea or observation]. Unlike prior approaches that [describe what existing methods do], our method [describe what makes your approach different and better].

% --- Example: how to reference a figure in Introduction ---
An overview of the proposed framework is illustrated in \cref{fig:method_overview}.

The main contributions of this paper are summarized as follows:
\begin{itemize}[leftmargin=2em,itemsep=2pt,topsep=2pt]
  \item We propose \method, a novel framework for [your task].
  \item We introduce [your key technique or insight], which enables [describe the benefit].
  \item Extensive experiments on [your benchmarks] demonstrate that \method achieves state-of-the-art performance, outperforming existing methods by [X\%] on [metric].
\end{itemize}
